\item \textbf{ACTIVIDAD} Está trabajando como asistente de enseñanza de un profesor de física. Él le proporcionó los siguientes datos sobre puntajes de exámenes para dos secciones diferentes de su clase de física:

\begin{center}
\begin{tabular}{|c|c|}
\hline
\textbf{Puntaje de exámenes Sección 1} & \textbf{Puntaje de exámenes Sección 2} \\ \hline
65 & 99 \\
65 & 95 \\
65 & 90 \\
65 & 85 \\
65 & 77 \\
65 & 75 \\
65 & 75 \\
65 & 73 \\
65 & 70 \\
65 & 67 \\
64 & 66 \\
64 & 65 \\
64 & 63 \\
64 & 58 \\
64 & 52 \\
64 & 49 \\
64 & 48 \\
64 & 35 \\
64 & 28 \\
64 & 20 \\ \hline
\end{tabular}
\end{center}

(a) Calcule el puntaje promedio para cada sección y el puntaje promedio rms para cada sección. (b) ¿Cómo se comparan los puntajes promedio de las dos secciones? (c) Compare los puntajes promedio y rms promedio de la Sección 1. ¿Por qué cree que estas dos cantidades tienen esta relación? (d) Compare el puntaje promedio y el puntaje promedio rms para la Sección 2. (e) ¿El promedio rms de un conjunto de números siempre será mayor que el promedio? ¿Por qué sí o por qué no?
